\ifdefined\gkver\else\def\gkver{student}\fi
\documentclass[\gkver]{gaokaozhenti}
\begin{document}

\chapter{2016年全国理综Ⅲ}
%% paperType: 理综物理部分

\begin{enumerate}
\item
%% number: 14
%% paperName: 2016年全国理综Ⅲ第 14 题 6 分
%% typeId: 1
%% type: 单选题
%% chapter: 其他
%% point: 物理学史
%% method: 无
%% score: 6
%% degree: 850
%% duplicateId: 0
%% body:
关于行星运动的规律，下列说法符合史实的是\xzanswer{B}
\fourchoices[answer=B]
{开普勒在牛顿定律的基础上，导出了行星运动的规律}
{开普勒在天文观测数据的基础上，总结出了行星运动的规律}
{开普勒总结出了行星运动的规律，找出了行星按照这些规律运动的原因}
{开普勒总结出了行星运动的规律，发现了万有引力定律}

%% answer: B
%% memo:
\memoanswer{
本题原卷及材料中未提供详解，待补充。
}

\item
%% number: 15
%% paperName: 2016年全国理综Ⅲ第 15 题 6 分
%% typeId: 1
%% type: 单选题
%% chapter: 电场
%% point: 等势面
%% method: 无
%% score: 6
%% degree: 750
%% duplicateId: 0
%% body:
关于静电场的等势面，下列说法正确的是\xzanswer{B}
\fourchoices[answer=B]
{两个电势不同的等势面可能相交}
{电场线与等势面处处相互垂直}
{同一等势面上各点电场强度一定相等}
{将一负的试探电荷从电势较高的等势面移至电势较低的等势面，电场力做正功}

%% answer: B
%% memo:
\memoanswer{
等势面相交，则电场线一定相交，故在同一点存在两个不同的电场强度方向，与事实不符，A 错误；电场线与等势面垂直，B 正确；同一等势面上的电势相同，但是电场强度不一定相同，C 错误；将一负电荷从高电势处移动到低电势处，受到的电场力的方向是从低电势指向高电势，所以电场力的方向与运动的方向相反，电场力做负功，D 错误。
}

\item
%% number: 16
%% paperName: 2016年全国理综Ⅲ第 16 题 6 分
%% typeId: 1
%% type: 单选题
%% chapter: 直线运动
%% point: 匀变速直线运动
%% method: 无
%% score: 6
%% degree: 750
%% duplicateId: 0
%% body:
一质点做速度逐渐增大的匀加速直线运动，在时间间隔 $t$ 内位移为 $s$，动能变为原来的9倍。该质点的加速度为\xzanswer{A}
\fourchoices[answer=A]
{$\frac{s}{t^{2}}$}
{$\frac{3s}{2t^{2}}$}
{$\frac{4s}{t^{2}}$}
{$\frac{8s}{t^{2}}$}

%% answer: A
%% memo:
\memoanswer{
本题原卷及材料中未提供详解，待补充。
}

\item
%% number: 17
%% paperName: 2016年全国理综Ⅲ第 17 题 6 分
%% typeId: 1
%% type: 单选题
%% chapter: 力物体的平衡
%% point: 三力平衡
%% method: 无
%% score: 6
%% degree: 750
%% duplicateId: 0
%% body:
如图，两个轻环a和b套在位于竖直面内的一段固定圆弧上：一细线穿过两轻环，其两端各系一质量为 $m$ 的小球。在a和b之间的细线上悬挂一小物块。平衡时，a、b间的距离恰好等于圆弧的半径。不计所有摩擦。小物块的质量为\xzanswer{C}
\onepicture[w=5.5cm]{figs/17.png}
\fourchoices[answer=C]
{$\frac{m}{2}$}
{$\frac{\sqrt{3}}{2}m$}
{$m$}
{$2m$}

%% answer: C
%% memo:
\memoanswer{
根据题意设悬挂小物块的点为 $O'$，圆弧的圆心为 $O$，由于 $ab=R$，所以三角形 $Oab$ 为等边三角形，根据几何知识可得 $\angle aO'b=120^{\circ}$，而一条绳子上的拉力相等，故 $T=mg$，小物块受到两条绳子的拉力作用且两力大小相等，夹角为 $120^{\circ}$，故受到的拉力的合力等于 $mg$，因为小物块受到绳子的拉力和重力作用，处于静止状态，故拉力的合力等于小物块的重力，为 $mg$，所以小物块的质量为 $m$，C 正确。
}

\item
%% number: 18
%% paperName: 2016年全国理综Ⅲ第 18 题 6 分
%% typeId: 1
%% type: 单选题
%% chapter: 磁场
%% point: 带电粒子在磁场中的运动
%% method: 无
%% score: 6
%% degree: 450
%% duplicateId: 0
%% body:
平面OM和平面ON之间的夹角为 $30^{\circ}$，其横截面（纸面）如图所示，平面OM上方存在匀强磁场，磁感应强度大小为 $B$，方向垂直于纸面向外。一带电粒子的质量为 $m$，电荷量为 $q$（$q>0$）。粒子沿纸面以大小为 $v$ 的速度从PM的某点向左上方射入磁场，速度与OM成 $30^{\circ}$ 角。已知粒子在磁场中的运动轨迹与ON只有一个交点，并从OM上另一点射出磁场。不计重力。粒子离开磁场的射点到两平面交线O的距离为\xzanswer{D}
\onepicture[w=4.8cm]{figs/18.png}
\fourchoices[answer=D]
{$\frac{mv}{2Bq}$}
{$\frac{\sqrt{3}mv}{Bq}$}
{$\frac{2mv}{Bq}$}
{$\frac{4mv}{Bq}$}

%% answer: D
%% memo:
\memoanswer{
本题原卷及材料中未提供详解，待补充。
}

\item
%% number: 19
%% paperName: 2016年全国理综Ⅲ第 19 题 6 分
%% typeId: 2
%% type: 多选题
%% chapter: 交流电
%% point: 变压器
%% method: 无
%% score: 6
%% degree: 550
%% duplicateId: 0
%% body:
如图，理想变压器原、副线圈分别接有额定电压相同的灯泡a和b。当输入电压 $U$ 为灯泡额定电压的10倍时，两灯泡均能正常发光。下列说法正确的是\xzanswer{AD}
\onepicture[w=4.2cm]{figs/19.png}
\fourchoices[answer=AD]
{原、副线圈匝数之比为9:1}
{原、副线圈匝数之比为1:9}
{此时 a 和 b 的电功率之比为9:1}
{此时 a 和 b 的电功率之比为1:9}

%% answer: AD
%% memo:
\memoanswer{
设灯泡的额定电压为 $U_{0}$，两灯泡均能正常发光，所以原线圈输出端电压为 $U_{1}=9U_{0}$，副线圈两端电压为 $U_{2}=U_{0}$，故 $\frac{U_{1}}{U_{2}}=\frac{9}{1}$，根据 $\frac{U_{1}}{U_{2}}=\frac{n_{1}}{n_{2}}=\frac{9}{1}$，A 正确，B 错误；根据公式 $\frac{I_{1}}{I_{2}}=\frac{n_{2}}{n_{1}}$ 可得 $\frac{I_{1}}{I_{2}}=\frac{1}{9}$，由于小灯泡两端的电压相等，所以根据公式 $P=UI$ 可得两者的电功率之比为1:9，C 错误，D 正确。
}

\item
%% number: 20
%% paperName: 2016年全国理综Ⅲ第 20 题 6 分
%% typeId: 2
%% type: 多选题
%% chapter: 曲线运动
%% point: 竖直平面圆周运动
%% method: 无
%% score: 6
%% degree: 450
%% duplicateId: 0
%% body:
如图，一固定容器的内壁是半径为 $R$ 的半球面；在半球面水平直径的一端有一质量为 $m$ 的质点 $P$。它在容器内壁由静止下滑到最低点的过程中，克服摩擦力做的功为 $W$。重力加速度大小为 $g$。设质点 $P$ 在最低点时，向心加速度的大小为 $a$，容器对它的支持力大小为 $N$，则\xzanswer{AC}
\onepicture[w=5.2cm]{figs/20.png}
\fourchoices[answer=AC]
{$a=\frac{2(mgR-W)}{mR}$}
{$a=\frac{2mgR-W}{mR}$}
{$N=\frac{3mgR-2W}{R}$}
{$N=\frac{2(mgR-W)}{R}$}

%% answer: AC
%% memo:
\memoanswer{
本题原卷及材料中未提供详解，待补充。
}

\item
%% number: 21
%% paperName: 2016年全国理综Ⅲ第 21 题 6 分
%% typeId: 2
%% type: 多选题
%% chapter: 电磁感应
%% point: 切割磁感线电动势
%% method: 无
%% score: 6
%% degree: 350
%% duplicateId: 0
%% body:
如图，M为半圆形导线框，圆心为 $O_{M}$；N是圆心角为直角的扇形导线框，圆心为 $O_{N}$；两导线框在同一竖直面（纸面）内；两圆弧半径相等；过直线 $O_{M}O_{N}$ 的水平面上方有一匀强磁场，磁场方向垂直于纸面。现使线框M、N在 $t=0$ 时从图示位置开始，分别绕垂直于纸面、且过 $O_{M}$ 和 $O_{N}$ 的轴，以相同的周期 $T$ 逆时针匀速转动，则\xzanswer{BC}
\onepicture[w=6.8cm]{figs/21.png}
\fourchoices[answer=BC]
{两导线框中均会产生正弦交流电}
{两导线框中感应电流的周期都等于 $T$}
{在 $t=\frac{T}{8}$ 时，两导线框中产生的感应电动势相等}
{两导线框的电阻相等时，两导线框中感应电流的有效值也相等}

%% answer: BC
%% memo:
\memoanswer{
本题原卷及材料中未提供详解，待补充。
}

\item
%% number: 22
%% paperName: 2016年全国理综Ⅲ第 22 题 5 分
%% typeId: 4
%% type: 实验题
%% chapter: 磁场
%% point: 安培力
%% method: 无
%% score: 5
%% degree: 550
%% duplicateId: 0
%% body:
某同学用图中所给器材进行与安培力有关的实验。两根金属导轨ab和 $a_{1}b_{1}$ 固定在同一水平面内且相互平行，足够大的电磁铁（未画出）的N极位于两导轨的正上方，S极位于两导轨的正下方，一金属棒置于导轨上且两导轨垂直。
\begin{enumerate}
	\item
	在图中画出连线，完成实验电路。要求滑动变阻器以限流方式接入电路，且在开关闭合后，金属棒沿箭头所示的方向移动。
	\item
	为使金属棒在离开导轨时具有更大的速度，有人提出以下建议：
	\threechoices[answer=AC]
	{适当增加两导轨间的距离}
	{换一根更长的金属棒}
	{适当增大金属棒中的电流}
	其中正确的是\tkanswer[1.2cm]{AC}（填入正确选项前的标号）
\end{enumerate}
\onepicture[w=8.5cm, align=right]{figs/22.png}

%% answer: 
\jdanswer{
\begin{enumerate}
	\item
	如图所示。
	\item
	AC
\end{enumerate}
}
%% memo:
\memoanswer{
本题原卷及材料中未提供详解，待补充。
}

\item
%% number: 23
%% paperName: 2016年全国理综Ⅲ第 23 题 10 分
%% typeId: 4
%% type: 实验题
%% chapter: 运动和力
%% point: 验证牛顿第二定律
%% method: 无
%% score: 10
%% degree: 350
%% duplicateId: 0
%% body:
某物理课外小组利用图（a）中的装置探究物体加速度与其所受合外力之间的关系。途中，置于试验台上的长木板水平放置，其右端固定一轻滑轮：轻绳跨过滑轮，一段与放在木板上的小滑车相连，另一端可悬挂钩码。本实验中可用的钩码共有 $N=5$ 个，每个质量均为 $0.010\Ukg$。实验步骤如下：
\begin{enumerate}
	\item
	将5个钩码全部放入小车中，在长木板左下方垫上适当厚度的小物快，使小车9（和钩码）可以在木板上匀速下滑。
	\item
	将 $n$（依次取 $n=1$，2，3，4，5）个钩码挂在轻绳右端，其余 $N-n$ 各钩码仍留在小车内；用手按住小车并使轻绳与木板平行。释放小车，同时用传感器记录小车在时刻 $t$ 相对于其起始位置的位移 $s$，绘制 $s-t$ 图像，经数据处理后可得到相应的加速度 $a$。
	\item
	对应于不同的 $n$ 的 $a$ 值见下表。$n=2$ 时的 $s-t$ 图像如图（b）所示：由图（b）求出此时小车的加速度（保留2位有效数字），将结果填入下表。
	\begin{center}
	\begin{tabular}{|c|c|c|c|c|c|}
	\hline
	$n$ & 1 & 2 & 3 & 4 & 5 \\
	\hline
	$a/\Umsq$ & 0.20 &  & 0.58 & 0.78 & 1.00 \\
	\hline
	\end{tabular}
	\end{center}
	\item
	利用表中的数据在图（c）中补齐数据点，并作出 $a-n$ 图像。从图像可以看出：当物体质量一定时，物体的加速度与其所受的合外力成正比。
	\item
	利用 $a-n$ 图像求得小车（空载）的质量为\tkanswer[1.2cm]{0.44}$\Ukg$（保留2位有效数字，重力加速度取 $g=9.8\Umsq$）。
	\item
	若以“保持木板水平”来代替步骤（1），下列说法正确的是\tkanswer[1.2cm]{B}（填入正确选项钱的标号）
	\threechoices[answer=B]
	{$a-n$ 图线不再是直线}
	{$a-n$ 图线仍是直线，但该直线不过原点}
	{$a-n$ 图线仍是直线，但该直线的斜率变大}
\end{enumerate}
\onepicture[w=6.0cm, align=right]{figs/23a.png}
\onepicture[w=10.0cm, align=right]{figs/23b.png}

%% answer: 
\jdanswer{
\begin{enumerate}
	\item
	——
	\item
	——
	\item
	0.40
	\item
	如图（c）所示。
	\item
	0.44
	\item
	B
\end{enumerate}
}
%% memo:
\memoanswer{
本题原卷及材料中未提供详解，待补充。
}

\item
%% number: 24
%% paperName: 2016年全国理综Ⅲ第 24 题 12 分
%% typeId: 6
%% type: 计算题
%% chapter: 曲线运动
%% point: 竖直平面圆周运动
%% method: 无
%% score: 12
%% degree: 550
%% duplicateId: 0
%% body:
如图，在竖直平面内由 $\frac{1}{4}$ 圆弧AB和 $\frac{1}{2}$ 圆弧BC组成的光滑固定轨道，两者在最低点B平滑连接。AB弧的半径为 $R$，BC弧的半径为 $\frac{R}{2}$。一小球在A点正上方与A相距 $\frac{R}{4}$ 处由静止开始自由下落，经A点沿圆弧轨道运动。
\begin{enumerate}
	\item
	求小球在 B、A 两点的动能之比；
	\item
	通过计算判断小球能否沿轨道运动到 C 点。
\end{enumerate}
\onepicture[w=4.2cm, align=right]{figs/24.png}

%% answer: 
\jdanswer{
\begin{enumerate}
	\item
	5
	\item
	小球恰好可以沿轨道运动到 C 点
\end{enumerate}
}
%% memo:
\memoanswer{
（1）设小球的质量为 $m$，小球在A点的动能为 $E_{kA}$，由机械能守恒可得 $E_{kA}=mg\frac{R}{4}$①。

设小球在B点的动能为 $E_{kB}$，同理有 $E_{kB}=mg\frac{5R}{4}$②。

由①②联立可得 $\frac{E_{kB}}{E_{kA}}=5$③。

（2）若小球能沿轨道运动到C点，小球在C点所受轨道的正压力 $N$ 应满足 $N\geq0$④。

设小球在C点的速度大小为 $v_{C}$，根据牛顿运动定律和向心加速度公式有 $N+mg=m\frac{v_{C}^{2}}{\frac{R}{2}}$⑤。

联立④⑤可得 $m\frac{2v_{C}^{2}}{R}\geq mg$⑥。

根据机械能守恒可得 $mg\frac{R}{4}=\frac{1}{2}mv_{C}^{2}$⑦。

根据⑥⑦可知，小球恰好可以沿轨道运动到 C 点。
}

\item
%% number: 25
%% paperName: 2016年全国理综Ⅲ第 25 题 20 分
%% typeId: 6
%% type: 计算题
%% chapter: 电磁感应
%% point: 电磁感应的电量
%% method: 无
%% score: 20
%% degree: 350
%% duplicateId: 0
%% body:
如图，两条相距 $l$ 的光滑平行金属导轨位于同一水平面（纸面）内，其左端接一阻值为 $R$ 的电阻；一与导轨垂直的金属棒置于两导轨上；在电阻、导轨和金属棒中间有一面积为 $S$ 的区域，区域中存在垂直于纸面向里的均匀磁场，磁感应强度打下 $B_{1}$ 随时间 $t$ 的变化关系为 $B_{1}=kt$，式中 $k$ 为常量；在金属棒右侧还有一匀强磁场区域，区域左边界MN（虚线）与导轨垂直，磁场的磁感应强度大小为 $B_{0}$，方向也垂直于纸面向里。某时刻，金属棒在一外加水平恒力的作用下从静止开始向右运动，在 $t_{0}$ 时刻恰好以速度 $v_{0}$ 越过MN，此后向右做匀速运动。金属棒与导轨始终相互垂直并接触良好，它们的电阻均忽略不计。求：
\begin{enumerate}
	\item
	在 $t=0$ 到 $t=t_{0}$ 时间间隔内，流过电阻的电荷量的绝对值；
	\item
	在时刻 $t$（$t>t_{0}$）穿过回路的总磁通量和金属棒所受外加水平恒力的大小。
\end{enumerate}
\onepicture[w=5.5cm, align=right]{figs/25.png}

%% answer: 
\jdanswer{
\begin{enumerate}
	\item
	$|q|=\frac{kt_{0}S}{R}$
	\item
	$\Delta\Phi_{t}=(B_{0}lv_{0}+kS)\Delta t$，$f=(B_{0}lv_{0}+kS)\frac{B_{0}l}{R}$
\end{enumerate}
}
%% memo:
\memoanswer{
本题原卷及材料中未提供详解，待补充。
}

\item
%% number: 33
%% paperName: 2016年全国理综Ⅲ第 33 题 10 分
%% typeId: 6
%% type: 计算题
%% chapter: 气体性质
%% point: 气体连接体
%% method: 无
%% score: 10
%% degree: 550
%% duplicateId: 0
%% body:
\begin{enumerate}
	\item
	（5 分）关于气体的内能，下列说法正确的是\xzanswer{CDE}。（填正确答案标号。选对1个得2分，选对2个得4分，选对3个得5分。每选错1个扣3分，最低得分为0分）
	\fivechoices[answer=CDE]
	{质量和温度都相同的气体，内能一定相同}
	{气体温度不变，整体运动速度越大，其内能越大}
	{气体被压缩时，内能可能不变}
	{一定量的某种理想气体的内能只与温度有关}
	{一定量的某种理想气体在等压膨胀过程中，内能一定增加}
	\item
	（10 分）一U形玻璃管竖直放置，左端开口，右端封闭，左端上部有一光滑的轻活塞。初始时，管内汞柱及空气柱长度如图所示。用力向下缓慢推活塞，直至管内两边汞柱高度相等时为止。求此时右侧管内气体的压强和活塞向下移动的距离。已知玻璃管的横截面积处处相同；在活塞向下移动的过程中，没有发生气体泄漏；大气压强 $p_{0}=75.0\UcmHg$。环境温度不变。
\end{enumerate}
\onepicture[w=5.8cm, align=right]{figs/33.png}

%% answer: 
\jdanswer{
\begin{enumerate}
	\item
	CDE
	\item
	$p_{1}'=144\UcmHg$，$h=9.42\Ucm$
\end{enumerate}
}
%% memo:
\memoanswer{
（2）设初始时，右管中空气柱的压强为 $p_{1}$，长度为 $l_{1}$；左管中空气柱的压强为 $p_{2}=p_{0}$，长度为 $l_{2}$。活塞被推下 $h$ 后，右管中空气柱的压强为 $p_{1}'$，长度为 $l_{1}'$；左管中空气柱的压强为 $p_{2}'$，长度为 $l_{2}'$。以 $\UcmHg$ 为压强单位。由题给条件得

$p_{1}=p_{0}+(20.0-5.00)\UcmHg$，

$l_{1}'=(20.0-\frac{20.0-5.00}{2})\Ucm$，

根据玻意耳定律得 $p_{1}l_{1}=p_{1}'l_{1}'$，

联立解得 $p_{1}'=144\UcmHg$。

根据题意可得 $p_{2}'=p_{1}'$，$l_{2}'=4.00\Ucm+\frac{20.0-5.00}{2}\Ucm-h$，

根据玻意耳定律得 $p_{2}l_{2}=p_{2}'l_{2}'$，解得 $h=9.42\Ucm$。
}

\item
%% number: 34
%% paperName: 2016年全国理综Ⅲ第 34 题 10 分
%% typeId: 6
%% type: 计算题
%% chapter: 光学
%% point: 光的折射
%% method: 无
%% score: 10
%% degree: 550
%% duplicateId: 0
%% body:
\begin{enumerate}
	\item
	（5 分）由波源S形成的简谐横波在均匀介质中向左、右传播。波源振动的频率为 $20\UHz$，波速为 $16\Ums$。已知介质中P、Q两质点位于波源S的两侧，且P、Q和S的平衡位置在一条直线上，P、Q的平衡位置到S的平衡位置之间的距离分别为 $15.8\Um$、$14.6\Um$，P、Q开始振动后，下列判断正确的是\xzanswer{BDE}。（填正确答案标号。选对1个得2分，选对2个得4分，选对3个得5分。每选错1个扣3分，最低得分为0分）
	\fivechoices[answer=BDE]
	{P、Q 两质点运动的方向始终相同}
	{P、Q 两质点运动的方向始终相反}
	{当 S 恰好通过平衡位置时，P、Q 两点也正好通过平衡位置}
	{当 S 恰好通过平衡位置向上运动时，P 在波峰}
	{当 S 恰好通过平衡位置向下运动时，Q 在波峰}
	\item
	（10 分）如图，玻璃球冠的折射率为 $\sqrt{3}$，其底面镀银，底面的半径是球半径的 $\frac{\sqrt{3}}{2}$ 倍；在过球心O且垂直于底面的平面（纸面）内，有一与底面垂直的光线射到玻璃球冠上的M点，该光线的延长线恰好过底面边缘上的A点。求该光线从球面射出的方向相对于其初始入射方向的偏角。
\end{enumerate}
\onepicture[w=3.2cm, align=right]{figs/34.png}

%% answer: 
\jdanswer{
\begin{enumerate}
	\item
	BDE
	\item
	$\beta=150^{\circ}$
\end{enumerate}
}
%% memo:
\memoanswer{
本题原卷及材料中未提供详解，待补充。
}

\item
%% number: 35
%% paperName: 2016年全国理综Ⅲ第 35 题 10 分
%% typeId: 6
%% type: 计算题
%% chapter: 运动和力
%% point: 动量和能量
%% method: 无
%% score: 10
%% degree: 450
%% duplicateId: 0
%% body:
\begin{enumerate}
	\item
	（5 分）一静止的铝原子核 \ce{^{27}_{13}Al} 俘获一速度为 $1.0\times10^{7}\Ums$ 的质子 p 后，变为处于激发态的硅原子核 \ce{^{28}_{14}Si}，下列说法正确的是\xzanswer{ABE}（填正确的答案标号，选对一个得2分，选对2个得4分，选对3个得5分，没错选1个扣3分，最低得分为零分）
	\fivechoices[answer=ABE]
	{核反应方程为 \ce{p + ^{27}_{13}Al -> ^{28}_{14}Si}}
	{核反应方程过程中系统动量守恒}
	{核反应过程中系统能量不守恒}
	{核反应前后核子数相等，所以生成物的质量等于反应物的质量之和}
	{硅原子核速度的数量级为 $10^{5}\Ums$，方向与质子初速度方向一致}
	\item
	（10 分）如图所示，水平地面上有两个静止的小物块a和b，其连线与墙垂直：a和b相距 $l$；b与墙之间也相距 $l$；a的质量为 $m$，b的质量为 $\frac{3}{4}m$，两物块与地面间的动摩擦因数均相同，现使a以初速度 $v_{0}$ 向右滑动，此后a与b发生弹性碰撞，但b没有与墙发生碰撞，重力加速度大小为 $g$，求物块与地面间的动摩擦因数满足的条件。
\end{enumerate}
\onepicture[w=6.5cm, align=right]{figs/35.png}

%% answer: 
\jdanswer{
\begin{enumerate}
	\item
	ABE
	\item
	$\frac{v_{0}^{2}}{2gl}\geq\mu\geq\frac{32v_{0}^{2}}{113gl}$
\end{enumerate}
}
%% memo:
\memoanswer{
根据质量数和电荷数守恒可得核反应方程 \ce{p + ^{27}_{13}Al -> ^{28}_{14}Si^*}，A 正确；过程中释放的核力远远大于外力，故系统动量守恒，B 正确；核反应过程中系统能量守恒，C 错误；由于反应过程中，要释放大量的能量，即伴随着质量亏损，所以生成物的质量小于反应物的质量之和，D 错误；由动量守恒可知，$mv=28mv'$，解得 $v'=\frac{1.0}{28}\times10^{7}\Ums$，故数量级约为 $10^{5}\Ums$，故 E 正确。
}

\end{enumerate}

\end{document}
